%% ====================================================================
\section{Conclusion}
\label{sec:bench-conclusion}
%% ====================================================================

The Mar-3 freeze ships a production-ready hybrid PQ consensus stack with measured numbers across every layer.

\paragraph{Headline numbers (Apple M1 Max single core, $n = 21$).}
\begin{itemize}[leftmargin=2em,nosep]
  \item BLS aggregate verify: $\sim 875$ \textmu s.
  \item ML-DSA-65 single verify: $\sim 250$ \textmu s; cached: $\sim 3$ \textmu s.
  \item Pulsar Pulse verify: $\sim 9$ ms.
  \item Full Prism Horizon verify (BLS + ML-DSA cert set + Pulse): $\sim 15$ ms.
  \item Warp 2.0 envelope verify (full Horizon, destination-side): $\sim 15$ ms.
  \item Cross-WAN $K = 21$ Pulse latency: $\sim 1$--$1.3$ s p99.
  \item LSS-Pulsar reshare $21 \to 21$: $\sim 960$ ms local-loop.
  \item Activation cert at $K = 21$ (cross-WAN, Lux mainnet): $\sim 880$ ms.
\end{itemize}

\paragraph{Operational reading.} The Lux mainnet bundle interval is $\sim 3$ s; the full Horizon path completes at $\sim 1$--$1.3$ s p99 cross-WAN, leaving $\sim 1.7$--$2$ s of headroom per bundle. The chain reaches Horizon-final on every bundle in steady state; Pulse-sealed lags Beam-final by $\sim 0.5$--$1$ s.

\paragraph{Where the headroom is.} Pulse verify ($\sim 9$ ms) is the dominant single cost; the GPU path (Section~\ref{sec:gpu-headroom}) projects 2--6$\times$ speedup at $K \in \{21, 64, 128\}$. ML-DSA cert-set verify ($\sim 5.3$ ms parallelized) reduces to $\sim 1$--$3$ ms with the Groth16 wrapper (Section~\ref{sec:prism-composition}).

\paragraph{What gates production GPU.} Constant-time review of GPU code paths; cross-platform KAT byte-equality (Metal + CUDA + WGSL); LP-137 \cite{lp137} milestones. The freeze tag ships CPU-only as the conservative production choice.

\paragraph{What scales with $n$.} ML-DSA cert-set verify is $\mathcal{O}(n)$ (per-validator); the Groth16 wrapper makes it $\mathcal{O}(1)$ at the cost of being classical. Pulse verify and Beam aggregate verify are $\mathcal{O}(1)$. For $n > 100$ deployments, grouped Pulsar (Section~\ref{sec:grouped-pulsar} of \cite{luxquasarhorizon}) keeps the per-group $K$ in the comfortable cross-WAN range.

\paragraph{What does not scale.} The Pulsar GroupKey ($\mathbf{A}$, $\tilde{\mathbf{b}}$ bytes) is byte-identical across reshare within an era. The PQ root of trust is one Pulsar verify regardless of $n$. The chain's PQ-floor latency is bounded by Pulse latency, not by $n$.

\paragraph{Cross-chain delivery.} Lux source $\to$ Hanzo destination on production mainnet: source Horizon-final $\sim 1.5$--$2$ s, Warp 2.0 envelope build $\sim 1$--$2$ ms, WAN delivery $\sim 100$--$300$ ms, destination verify $\sim 15$ ms. Total source-event $\to$ destination-final: $\sim 2$--$3$ s. PQ source-finality, no oracle, no separate guardian network.

\paragraph{What is honest, what is wrapper.} Pulse + ML-DSA cert set are the PQ root of trust. The Groth16 wrapper around the ML-DSA cert set is a classical compression artifact; not a PQ contribution. Numbers are reported with the wrapper status flagged (Section~\ref{sec:prism-composition}); no claim in this report obscures the classical-vs-PQ distinction.

\paragraph{Reproducibility.} Every table is reproducible via the listed command from the freeze tag. KAT regen is byte-deterministic. The fuzz harnesses return zero panics. The CI gate is \texttt{GOWORK=off go test -count=1 -short -timeout 300s ./...} from the consensus, pulsar, lens, threshold, and warp roots.

The Mar-3 freeze ships the production numbers. The benchmark report establishes the baseline; subsequent improvements (GPU production, grouped-Pulsar production at $G > 1$, Lumen transport) move from this baseline.
